\section{A complete criterion for centered harmonic pole cancellation}
\label{hc:section}

Question~11 of the incoming Exact Resonant Identities report \cite{RepoExact} asks which
polynomials in generalized harmonic numbers have no principal part at
any nonpositive even spectral integer when the outer shift is $1/2$.
The criterion below is finite and necessary as well as sufficient.
The necessity is the substantive point: parity of the individual
asymptotic expansions alone supplies only one implication.

For a fixed $R\ge1$, write
\[
 H_n^{(r)}=\sum_{k=1}^n k^{-r},\qquad H_n^{(1)}=H_n,
 \qquad
 D_P(s)=\sum_{n=0}^{\infty}
 \frac{P(H_n^{(1)},\ldots,H_n^{(R)})}{(n+1/2)^s},
 \quad P\in\mathbb C[X_1,\ldots,X_R].
\]
The defining series is absolutely convergent for $\Re s>1$.
As usual $H_0^{(r)}=0$, so the finite $n=0$ term is retained.
Define an involution $\iota$ on the polynomial ring by
\begin{equation}\label{hc:involution}
 \iota(X_{2k})=2\zeta(2k)-X_{2k},\qquad
 \iota(X_{2k+1})=X_{2k+1},\qquad \iota(X_1)=X_1.
\end{equation}
Only variables whose indices are at most $R$ are used.

\begin{theorem}[Necessary and sufficient centered cancellation]
\label{hc:main}
The function $D_P$ has a meromorphic continuation to $\mathbb C$.
It is holomorphic at every point $s=0,-2,-4,\ldots$ if and only if
\begin{equation}\label{hc:criterion}
                         \iota(P)=P.
\end{equation}
Equivalently, put $Y_{2k}=X_{2k}-\zeta(2k)$. After expressing $P$ in
the variables $X_1,X_3,X_5,\ldots$ and $Y_2,Y_4,Y_6,\ldots$, every
monomial that occurs must have even total degree in the $Y$ variables.
The invariant algebra is therefore
\begin{equation}\label{hc:invariant-ring}
 \mathbb C\bigl[X_1,X_3,X_5,\ldots,
        Y_{2i}Y_{2j}\ (1\le i\le j\le\lfloor R/2\rfloor)\bigr].
\end{equation}
This is a generating description, not a claim that the displayed
quadratic generators are algebraically independent.
\end{theorem}

For example, every polynomial in odd-order harmonic numbers has the
cancellation, as does every product of an even number of even-order
tails. Thus the numerators
\[
 H_n^{(3)},\qquad
 H_n^2H_n^{(5)},\qquad
 (H_n^{(2)}-\zeta(2))^2,\qquad
 (H_n^{(2)}-\zeta(2))(H_n^{(4)}-\zeta(4))
\]
all satisfy the criterion. The two even-order factors need not be equal.
The raw numerator $(H_n^{(2)})^2$ fails: its simple-pole coefficient
at $s=0$ is $-2\zeta(2)$.

\subsection{Bernoulli coefficients and the principal parts}

Use $x=n+1/2$, $z=x^{-1}$, and the formal symbol $\mathcal L=\log x$. Define
\begin{align}
 h_1(z,\mathcal L)&=\mathcal L+\gamma-
           \sum_{k\ge1}\frac{B_{2k}(1/2)}{2k}z^{2k},\label{hc:h1}\\
 h_r(z)&=\zeta(r)-\frac{z^{r-1}}{r-1}
  -\sum_{k\ge1}\frac{B_{2k}(1/2)}{(2k)!}(r)_{2k-1}
                         z^{r+2k-1},\qquad r\ge2.\label{hc:hr}
\end{align}
These are formal asymptotic series; no convergence of the infinite
Bernoulli expansions is asserted. Each finite truncation is the
usual shifted polygamma or Hurwitz-zeta asymptotic expansion \cite{DLMFGamma,DLMFHurwitz}, with
a remainder sufficient for termwise Dirichlet subtraction.
Expand
\begin{equation}\label{hc:A}
 P(h_1,h_2,\ldots,h_R)=\sum_{j\ge0}A_j(\mathcal L)z^j,
             \qquad A_j\in\mathbb C[\mathcal L].
\end{equation}
Every fixed $A_j$ is a finite calculation. The degree in $\mathcal L$ is
bounded by $\deg_{X_1}P$.

\begin{proposition}[All principal coefficients]
\label{hc:principal}
At $s=1-j$, $j\ge0$, the entire Laurent principal part of $D_P$ is
\begin{equation}\label{hc:principal-formula}
 \sum_{\ell\ge0}
 \frac{\ell!\,[\mathcal L^\ell]A_j(\mathcal L)}{(s+j-1)^{\ell+1}}.
\end{equation}
Consequently $D_P$ is regular at all nonpositive even integers
exactly when $A_{2m+1}=0$ for every $m\ge0$.
\end{proposition}
\begin{proof}
Subtract finitely many terms of \eqref{hc:A} from the defining
series. Each subtracted term sums to
$(-1)^\ell\partial_s^\ell\zeta(s+j,1/2)$ times its coefficient.
The finite asymptotic remainder gives a normally convergent series
in a successively larger left half-plane. These formulas agree on
overlaps and continue $D_P$ meromorphically. The only principal
coefficient of $(-1)^\ell\partial_s^\ell\zeta(s+j,1/2)$ is
$\ell!/(s+j-1)^{\ell+1}$, proving the claim. Distinct logarithmic
degrees give distinct Laurent powers, so none is lost by merely
checking the residue.
\end{proof}

Here are the first three finite Bernoulli tests. Let $\partial_r$
denote polynomial differentiation with respect to $X_r$, omitting
indices larger than $R$, and define commuting operators
\begin{align*}
 A&=-\partial_2,&
 B&=\frac{\partial_1}{24}-\frac{\partial_3}{2},&
 C&=\frac{\partial_2}{12}-\frac{\partial_4}{3},\\
 D&=-\frac{7\partial_1}{960}+\frac{\partial_3}{8}
                    -\frac{\partial_5}{4},&
 E&=-\frac{7\partial_2}{240}+\frac{\partial_4}{6}
                    -\frac{\partial_6}{5}.
\end{align*}
Evaluate the following expressions at
$\mathbf b(\mathcal L)=(\mathcal L+\gamma,\zeta(2),\ldots,\zeta(R))$:
\begin{align}
 A_1(\mathcal L)&=(AP)(\mathbf b(\mathcal L)),\notag\\
 A_3(\mathcal L)&=\left((C+AB+A^3/6)P\right)(\mathbf b(\mathcal L)),\label{hc:first-tests}\\
 A_5(\mathcal L)&=\left((E+AD+BC+(A^2C+AB^2)/2
                   +A^3B/6+A^5/120)P\right)(\mathbf b(\mathcal L)).\notag
\end{align}
They follow by expanding the finite translation operator
$\exp(zA+z^2B+z^3C+z^4D+z^5E+\cdots)$.
For example,
\[
 A_3=\left(\frac{P_2}{12}-\frac{P_4}{3}
   -\frac{P_{12}}{24}+\frac{P_{23}}2-\frac{P_{222}}6\right)(\mathbf b(\mathcal L)).
\]
The rational-coefficient polynomial
$P=X_2^3-\tfrac{15}{2}X_2X_4$ satisfies $A_1=0$, because
$\zeta(4)=2\zeta(2)^2/5$, but
$A_3=\tfrac52\zeta(2)-1\ne0$. Thus its Dirichlet series is regular
at zero and has a simple pole at $-2$. More generally
$P=X_{2M+2}-\zeta(2M+2)$ passes all earlier tests and first has a
pole at $s=-2M$, with residue $-1/(2M+1)$. Finite lists of tested
spectral centers cannot replace the polynomial criterion when $R$
is unrestricted.

\subsection{Why the parity criterion is necessary}

The following elementary difference-algebra lemma is the relevant
form of the classical additive-difference mechanism behind
H\"older's theorem; compare the general theory in \cite{HardouinSinger}. It is proved here, so the formal-series application
does not rely on a converse assertion about analytic asymptotic
expansions.

\begin{lemma}[A polynomial relation forces rational telescoping]
\label{hc:additive}
Let $K=\mathbb C(x)$ with $\sigma x=x+1$, and let $E$ be a field
extension carrying an extension of $\sigma$. Suppose
$f_1,\ldots,f_d\in E$ satisfy $\sigma f_i=f_i+a_i$, $a_i\in K$.
If the $f_i$ are algebraically dependent over $K$, then there exist
constants $c_i\in\mathbb C$, not all zero, and $b\in K$ such that
\begin{equation}\label{hc:telescoping}
                       \sum_i c_i a_i=\sigma b-b.
\end{equation}
\end{lemma}
\begin{proof}
Choose a nonzero polynomial relation $Q(\mathbf f)=0$ of least
total degree $d_0$. Among such relations choose one with the fewest
nonzero monomials of degree $d_0$, and divide by one of its leading
coefficients. That coefficient is now one. The polynomial
\[
              \sigma Q(\mathbf X+\mathbf a)-Q(\mathbf X)
\]
also vanishes at $\mathbf f$. Its chosen leading monomial cancels,
and no new degree-$d_0$ monomial is created. Minimality first of
degree and then of leading support forces this difference to be
the zero polynomial. If $Q_{d_0}$ is the highest homogeneous part,
then $\sigma Q_{d_0}=Q_{d_0}$. Every rational periodic coefficient
is constant, so $Q_{d_0}\in\mathbb C[\mathbf X]$.
Comparison in degree $d_0-1$ gives
\[
 \sigma Q_{d_0-1}-Q_{d_0-1}
              =-\sum_i a_i\,\partial_i Q_{d_0}.
\]
At least one coefficient of the gradient on the right is a
nonzero vector in $\mathbb C^d$, since $d_0\ge1$ and the
characteristic is zero. Taking that monomial coefficient and
reversing the sign of the rational coefficient on the left proves
\eqref{hc:telescoping}.
\end{proof}

\begin{lemma}[Independence of the formal centered harmonic series]
\label{hc:independence}
For every finite $R$, the formal series $h_1,\ldots,h_R$ in
\eqref{hc:h1}--\eqref{hc:hr} are algebraically independent over
$\mathbb C(x)$, where $x=z^{-1}$ and $\mathcal L$ is a formal logarithm.
\end{lemma}
\begin{proof}
Work in $E=\mathbb C((z))(\mathcal L)$. Extend translation by
\[
 \sigma z=\frac{z}{1+z},\qquad
 \sigma \mathcal L=\mathcal L+\log(1+z),
\]
and use the derivation $\mathcal D z=-z^2$, $\mathcal D \mathcal L=z$.
These operations commute. Put $f=h_1-\gamma$.
The Bernoulli expansion and the digamma recurrence give the formal
identity
\[
 \sigma f-f=\frac1{x+1/2},\qquad
 \sigma\mathcal D^j f-\mathcal D^j f
       =\frac{(-1)^j j!}{(x+1/2)^{j+1}}\quad(j\ge0).
\]
It can equally be checked coefficientwise from the Bernoulli
generating function. A nonzero constant linear combination of the
right sides has exactly one finite pole, at $-1/2$.

No nonzero rational difference $b(x+1)-b(x)$ has just one finite
pole. Indeed partition the finite poles of $b$ into classes modulo
$\mathbb Z$. In any nonempty class, the extreme poles of $b$ give
two distinct, uncanceled extreme poles of its difference. If $b$
has no finite poles, it is polynomial and its difference is
polynomial. Lemma~\ref{hc:additive} therefore proves algebraic
independence of $f,\mathcal Df,\ldots,\mathcal D^{R-1}f$.
Finally
\[
 h_1=f+\gamma,\qquad
 h_r=\zeta(r)-\frac{(-1)^r}{(r-1)!}\mathcal D^{r-1}f\quad(r\ge2)
\]
is an invertible affine change of these variables.
\end{proof}

\begin{proof}[Proof of Theorem~\ref{hc:main}]
Let $\tau$ act on formal expansions by $z\mapsto-z$, $\mathcal L\mapsto \mathcal L$.
Equations \eqref{hc:h1}--\eqref{hc:hr} show
$\tau h_r=\iota(h_r)$: the even harmonic tails change sign and
the odd harmonic variables do not. By
Proposition~\ref{hc:principal}, absence of every nonpositive even
principal part is equivalent to
$\tau P(\mathbf h)=P(\mathbf h)$, hence to
$(\iota P-P)(\mathbf h)=0$. Lemma~\ref{hc:independence} makes
this equivalent to the polynomial identity $\iota P=P$.
In centered variables $Y_{2k}$ the involution changes all signs
simultaneously. Its fixed monomials have even total $Y$-degree,
and each such monomial factors into quadratic pairs. This proves
\eqref{hc:invariant-ring}.
\end{proof}

\begin{corollary}[Algebraic coefficients]
\label{hc:algebraic}
If $P$ has algebraic-number coefficients, the criterion holds if
and only if $P$ is independent of every even-index variable.
\end{corollary}
\begin{proof}
Write $\zeta(2k)=c_k\pi^{2k}$ with nonzero $c_k\in\mathbb Q$.
Since $\pi^2$ is transcendental, invariance of a polynomial over
$\overline{\mathbb Q}$ under \eqref{hc:involution} implies its
invariance under the polynomial reflection
$X_{2k}\mapsto2c_kT^k-X_{2k}$ for every indeterminate $T$.
Composing two such reflections gives invariance under translations
by $2(c_kT^k-c_kU^k)_k$. Differentiating in $T$ and then setting
$T=U$ yields
\[
             \sum_k k c_k U^{k-1}\partial_{2k}P=0.
\]
Distinct powers of $U$ force each $\partial_{2k}P$ to vanish.
The converse follows directly from Theorem~\ref{hc:main}.
\end{proof}

This corollary uses only the established transcendence of $\pi$.
The main theorem makes no assertion about algebraic independence
of evaluated zeta or Stieltjes constants.

\subsection{All centered even moments of a trigamma square}

The invariant numerator $(H_n^{(2)}-\zeta(2))^2$ admits a particularly
short all-degree value formula. Put
\[
 T_n=\zeta(2,n+1)=\psi'(n+1),\qquad
 V_M=D_{(X_2-\zeta(2))^2}(-2M),\quad M\ge0.
\]
Every point in this definition is regular by Theorem~\ref{hc:main}.

\begin{theorem}[Finite Bernoulli formula for all even moments]
\label{hc:trigamma}
For every $M\ge0$,
\begin{align}\label{hc:all-moments}
 V_M={}&3B_{2M}(1/2)\zeta(3)\notag\\
 &-\frac1{2(2M+1)}\sum_{\ell=1}^{M}
 \binom{2M+1}{2M-2\ell}B_{2M-2\ell}(1/2)
                     (2\ell-1)B_{2\ell-2}.
\end{align}
The sum is empty at $M=0$. In particular all these values belong
to $\mathbb Q+\mathbb Q\zeta(3)$.
\end{theorem}
\begin{proof}
The polynomial partial sum is
\[
 S_M(k)=\sum_{n=0}^{k-1}(n+1/2)^{2M}
 =\frac{B_{2M+1}(k+1/2)}{2M+1}
 =\sum_{\ell=0}^M s_{M,\ell} k^{2\ell+1},
\]
where
$s_{M,\ell}=\binom{2M+1}{2M-2\ell}B_{2M-2\ell}(1/2)/(2M+1)$.
Since $T_{k-1}-T_k=k^{-2}$, finite summation by parts gives
\[
 \sum_{n=0}^{N-1}(n+1/2)^{2M}T_n^2
 =\sum_{\ell=0}^M s_{M,\ell}E_{2\ell+1}(N),
\]
with
\[
 E_q(N)=N^qT_N^2+2\sum_{k=1}^N k^{q-2}T_{k-1}
                                -\sum_{k=1}^N k^{q-4}.
\]
At $q=1$ the limit is $3\zeta(3)$: the classical Euler sum \cite{Can}
$\sum_{k\ge1}H_k/k^2=2\zeta(3)$ implies
$\sum_{k\ge1}T_{k-1}/k=2\zeta(3)$, by interchanging a positive
double sum. The remaining diagonal term is $\zeta(3)$.

For odd $q\ge3$, let $P_{q-2}(N)=\sum_{k=1}^N k^{q-2}$.
Interchanging a finite sum gives
\[
 \sum_{k=1}^N k^{q-2}T_{k-1}
   =\sum_{j=1}^N\frac{P_{q-2}(j)}{j^2}
                                +P_{q-2}(N)T_N.
\]
In Faulhaber's polynomial the half-leading term cancels
$\sum k^{q-4}$ from $E_q$. All other finite power sums are
polynomials in $N$ with zero constant coefficient. Hence the
constant term of the expansion of $E_q(N)$ at infinity is the
constant of $N^qT_N^2+2P_{q-2}(N)T_N$.
Use
\[
 T_N=N^{-1}-\tfrac12N^{-2}
                   +\sum_{j\ge1}B_{2j}N^{-2j-1}.
\]
Only finitely many terms are needed. The first product contributes
$-B_{q-3}$ and the second contributes
$B_{q-3}-(q-2)B_{q-3}/2$. Their sum is
$-(q-2)B_{q-3}/2$.

It remains to identify this cutoff constant with the stated
spectral value. The centered expansion of $T_n^2$ contains only
even powers of $(n+1/2)^{-1}$. Subtract from
$(n+1/2)^{2M}T_n^2$ all its nondecaying even powers. The residual
is $O(n^{-2})$ and is ordinarily summable. Each subtracted power
has spectral value $\zeta(-2j,1/2)=0$, and its finite partial sum
is an odd polynomial in $N$ with zero constant. Thus the residual
sum, the cutoff constant in $N$, and $V_M$ coincide. Combining
the coefficients $s_{M,\ell}$ proves \eqref{hc:all-moments}.
\end{proof}

The first four identities can be written as ordinary convergent
bracketed series, with $x_n=n+1/2$:
\begin{align}
 \sum_{n\ge0}T_n^2&=3\zeta(3),\notag\\
 \sum_{n\ge0}\{x_n^2T_n^2-1\}
        &=-\frac16-\frac14\zeta(3),\notag\\
 \sum_{n\ge0}\{x_n^4T_n^2-x_n^2+\tfrac16\}
        &=\frac1{30}+\frac7{80}\zeta(3),\label{hc:moment-examples}\\
 \sum_{n\ge0}\{x_n^6T_n^2-x_n^4+\tfrac16x_n^2-\tfrac{47}{720}\}
        &=\frac1{672}-\frac{31}{448}\zeta(3).\notag
\end{align}
Each bracket is $O(n^{-2})$. These statements do not assign an
ordinary sum to any of the divergent unbracketed expressions.

